% Source: https://artofproblemsolving.com/community/c3h1563396_orthocenter
% Stable identity: aops-topic-1563396-q01
Let $O$ be the origin and $\vec u_i=(\cos t_i,\sin t_i)$, unit vectors with $\vec u_i\cdot\vec u_j=\cos(t_i-t_j)$. The points lie on the circle of radius $a$ about $O$, so $O$ is the circumcentre of $ABC$. Now
\[
\left|\vec u_1+\vec u_2+\vec u_3\right|^2=3+2\bigl(\cos(t_1-t_2)+\cos(t_2-t_3)+\cos(t_3-t_1)\bigr)=3+2\left(-\dfrac{3}{2}\right)=0 .
\]
So $\vec u_1+\vec u_2+\vec u_3=\vec0$, and the centroid $G=\dfrac {a}{3}\left(\vec u_1+\vec u_2+\vec u_3\right)$ is $O$: the centroid equals the circumcentre.

\begin{asy}
usepackage("amsmath");
size(5cm);
defaultpen(fontsize(10pt));
real a=1;
real[] t={20,140,260};
pair O=(0,0);
draw(circle(O,a),gray);
pair A=dir(t[0]),B=dir(t[1]),C=dir(t[2]);
draw(A--B--C--cycle,linewidth(0.9));
draw(O--A,blue,Arrow(TeXHead));draw(O--B,blue,Arrow(TeXHead));draw(O--C,blue,Arrow(TeXHead));
pair M=(B+C)/2;draw(O--M,dashed);
dot(O);label("$O$",O,1.6*dir(85));
label("$A$",A,NE);label("$B$",B,NW);label("$C$",C,S);label("$M$",M,W);
label("\mbox{}",(min(currentpicture,user=true)+max(currentpicture,user=true))/2,black); // keeps the text after the figure black when LaTeX typesets the labels inline
\end{asy}

\textbf{Centroid = circumcentre forces equilateral.} Let $M$ be the midpoint of $BC$. The centroid lies on the median $AM$ with $GM=\dfrac{1}{3}AM>0$, so $O=G\ne M$. The circumcentre lies on the perpendicular bisector of $BC$, which is the line through $M$ perpendicular to $BC$. The line $OM$ is therefore that perpendicular bisector, and it is also the median line $AM$ (it contains $A$, $G=O$ and $M$). So $A$ is on the perpendicular bisector of $BC$, and $AB=AC$. In the same way $BA=BC$. Hence $AB=BC=CA$ and the triangle is equilateral. $\blacksquare$
